% ============================================================================
%  Assumptions invoked by the corrected Proposition prp:lepskii
%  (Lepskii oracle inequality for the location-scale LF estimator)
%
%  These assumptions either restate the source-paper assumptions
%  (\ref{ass:identif} -- \ref{ass:restrvarth}; \ref{ass:app1} -- \ref{ass:app6})
%  or strengthen / replace them as required by the audit patches.
%  Labels are stable across results.tex / proofs.tex.
% ============================================================================

% ----------------------------------------------------------------------------
% Restated identification scaffolding (verbatim from the source paper)
% ----------------------------------------------------------------------------

\begin{assumption}[Local identification]\label{ass:identif}
$b_\dagger \in \mathcal{A}'_b$ and the model is locally identified on
$\mathcal{B}_{b_\dagger}$ in the sense of Proposition~\ref{prp:onetoone}.
\end{assumption}

\begin{assumption}[Initial condition]\label{ass:initcond}
$b_0 \in \mathcal{B}_{b_\dagger}$.
\end{assumption}

\begin{assumption}[Source condition]\label{ass:sourcecond}
Under Assumption~\ref{ass:initcond}, there exists $\beta \in (0, 2]$ such that
$b'_0 - b'_\dagger = (\mathsf{A}^*_{b_\dagger}\mathsf{A}_{b_\dagger})^{\beta/2} v$
for some $v \in \mathcal{H}$ with $\|v\| = 1$.
\end{assumption}

\begin{assumption}[Link condition]\label{ass:linkcond}
For some $a > 1$, $(\mathsf{T}'^{*}_{b_\dagger}\mathsf{T}'_{b_\dagger})^{1/(2a)}
\sim D^{-1}$ in the matched Hilbert-scale sense of \citet{florens2012}.
\end{assumption}

\begin{assumption}[Frechet differentiability of $\hat{\mathsf T}$]\label{ass:frediff}
\textup{(i)} $\hat{\mathsf T}$ is Frechet differentiable for all
$b \in \mathcal{A}'$ in the $b$-coordinate system.
% PATCH-R3-G2-a: C^2 Fréchet hypothesis on the Hessian-like operator.
\textup{(ii)} The map $b \mapsto \hat{\mathsf T}'^{*}_{b}\hat{\mathsf T}'_{b}$
is twice Fr\'echet differentiable on $\mathcal{B}_{b_\dagger}$, with first
Fr\'echet derivative denoted $\mathsf{L}_{b}$ (a bounded operator on the
ambient Hilbert space) and second Fr\'echet derivative bounded in
operator norm by a deterministic constant $L'_{\mathsf T} < \infty$ on
$\mathcal{B}_{b_\dagger}$, uniformly on the event $\Omega'_n$ of
Lemma~\ref{lem:propagator}. Equivalently, there exists
$L'_{\mathsf T} < \infty$ such that, for all $b, b' \in \mathcal{B}_{b_\dagger}$
on $\Omega'_n$,
\[
\bigl\|\hat{\mathsf T}'^{*}_{b}\hat{\mathsf T}'_{b}
- \hat{\mathsf T}'^{*}_{b'}\hat{\mathsf T}'_{b'}
- \mathsf{L}_{b'}(b - b')\bigr\|
\;\le\; L'_{\mathsf T}\,\|b - b'\|^{2}.
\]
% PATCH-R4-B5: disclose explicit dependence list for L'_T.
The constant
$L'_{\mathsf T} = L'_{\mathsf T}(R^\sigma, M_1, M_2, M_3,
\kappa_n^{-1}, \text{kernel-envelope}) < \infty$ on $\Omega'_n$.
The $\kappa_n^{-1}$-factor enters through the Hessian of
$\hat{\mathsf T}'^{*}_{b}\hat{\mathsf T}'_{b}$ along the
$f_U^{-1}$-direction; under
Assumption~\ref{ass:restrvarth}, the product
$\kappa_n^{-1}\delta_n^{2}M_n^{2}V(M_n)^{2} = o(1)$ at $m = M_n$, so
the propagator-residual absorption in Step~3 of
Proposition~\ref{prp:lepskii} closes.
\end{assumption}

% ----------------------------------------------------------------------------
% PATCH-M6 / PATCH-CR1 / PATCH-MN10:
%   Strengthened operator-rate assumption with $L^{2l}$ moments
%   (replaces source ass:rateconv).
% ----------------------------------------------------------------------------

\begin{assumption}[$L^{2l}$ rates of convergence of estimated operators]\label{ass:rateconv}
\textup{(i)} There exist sequences $\delta_n, \gamma_n \to 0$ as
$n \to \infty$ and an integer $l_0 \ge 4$ such that, for every integer $l$
with $1 \le l \le l_0$,
\[
E\,\bigl\|\hat{\mathsf T}(b_\dagger) - \mathsf T(b_\dagger)\bigr\|^{2l}
\;=\; O\!\bigl(\delta_n^{2l}\bigr),
\qquad
E\,\bigl\|\hat{\mathsf A}^{*}_{b_\dagger} - \mathsf A^{*}_{b_\dagger}\bigr\|^{2l}
\;=\; O\!\bigl(\gamma_n^{2l}\bigr).
\]
The rates $(\delta_n, \gamma_n)$ specialise to the discrete-$W$ (FI) or
continuous-$W$ (MI) regime as in the source paper.

% PATCH-R3-G3 (NEW): hypercontractive / moment-equivalence clause on iterates.
\textup{(iv)} \emph{(Moment equivalence on iterates.)} There exists a
finite constant $C_q > 0$ such that, for every integer $q \ge 1$ and every
$m \le M_n$,
\[
\bigl(E\,\|\hat e_m\|^{2q}\bigr)^{1/q}
\;\le\; C_q\, E\,\|\hat e_m\|^{2},
\qquad
\hat e_m := \hat b'_m - b'_\dagger,
\]
and the same inequality holds with $\hat e_m$ replaced by
$\mathsf T'_{b_\dagger}\hat e_m$, with constant $C_q$ that may differ but
is finite for every $q \le 2l_0$. Equivalently, the iterate $\hat e_m$ is
sub-Gaussian on $\mathcal{B}_{b_\dagger}$ with proxy variance
$C_1\,E\|\hat e_m\|^{2}$.
\end{assumption}

\begin{remark}[Justification of Assumption \ref{ass:rateconv}(iv)]\label{rem:hyper}
The moment-equivalence clause is plausible under the surrounding
assumptions but is imposed, not derived. Under
Assumption~\ref{ass:subgauss} plus the kernel-density-based construction
of $\hat e_m$, the four-term decomposition
$\hat e_m = I_m + II_m + III_m + IV_m$ used in
Lemma~\ref{lem:itervar} has each term sub-Gaussian individually:
Term~$II_m$ and Term~$IV_m$ are sub-Gaussian kernel sample-averages
weighted by the bounded LF spectral filter (Hansen, 2008,
Theorem~6); Term~$III_m$ is a quadratic Fr\'echet remainder
deterministically bounded on $\Omega_n$; Term~$I_m$ is initial-condition
bias built from sub-Gaussian pilots. Sums of sub-Gaussian variables are
sub-Gaussian, and the moment-equivalence
$(E|X|^{2q})^{1/q} \le C_q\,E|X|^2$ for a sub-Gaussian real random
variable $X$ is the standard hypercontractive principle of
\citet[Theorem~2.5.2]{vershynin2018}, with the foundational entropy
form in \citet[Ch.~2]{boucheron2013} as supporting context. We state
\textup{(iv)} as an assumption rather than derive it because the
precise sub-Gaussian proxy of the sum requires control of
cross-correlations between the four terms that is technical and not
informative for the broader Lepskii argument.
\end{remark}

\begin{assumption}[Tuning parameters]\label{ass:restrvarth}
With $\kappa_n = \inf_{|u|\le \ell_n} f_U(u)$ and $\epsilon \in (0, 1/2)$:
\textup{(i)} $\delta_n N = O(1)$ and $\gamma_n N^{(\epsilon+1)/2} = o(1)$;
\textup{(ii)} $(\delta_n \vee \gamma_n)\sqrt{N}/(h_u^2 \kappa_n) = o(1)$;
\textup{(iii)} $N^{-\epsilon/2}/(h_u^2 \kappa_n) = O(1)$.
\end{assumption}

% ----------------------------------------------------------------------------
% PATCH-M5: uniform-in-x bound on phi_dagger
% (this is the strengthened version of ass:unifbd from the source paper).
% ----------------------------------------------------------------------------

\begin{assumption}[Uniform bounds; strengthened]\label{ass:unifbd}
\textup{(i)} The conditional density $f_{Y\mid X,W}(y\mid x,w)$ is uniformly
bounded above by $M_1 < \infty$.
\textup{(ii)} $|\partial_y f_{Y\mid X,W}(y\mid x,w)| \le M_2$ and
$f_{X\mid W}(x\mid w) \le M_3 f_X(x)$.
\textup{(iii)} The scale and location functions are uniformly bounded:
$0 < R_\sigma \le \sigma_\dagger(x) \le R^\sigma < \infty$ and
$|\varphi_\dagger(x)| \le R^\varphi < \infty$ for all $x$. Equivalently,
$-\log R^\sigma \le \omega_\dagger(x) \le -\log R_\sigma$ and
$\|g_\dagger\|_\infty < \infty$.
\textup{(iv)} $E[U^4] \le m_4 < \infty$.
\end{assumption}

% Auxiliary regularity from the source paper (restated verbatim for
% completeness; cited as a block).

\begin{assumption}[Density and kernel regularity]\label{ass:app1to3}
The conditional density $f_{Y,X\mid W}$ and the marginal $f_{Y,X}$ are
$d$-times continuously differentiable and uniformly bounded; $f_{X,W}$, $f_X$,
$f_W$ are bounded away from $0$ and $\infty$; $f_U$ is absolutely continuous
and $d$-times continuously differentiable; the kernels $C_h$ (continuous $W$)
and $L_{h_w}$ (discrete $W$) satisfy the order-$\rho$ and Aitchison-Aitken
conditions of \textup{ass:app3a}--\textup{ass:app3b} in the source paper.
\end{assumption}

\begin{assumption}[Residual density floor]\label{ass:app4}
There exists a sequence $\ell_n$, either bounded or slowly diverging,
such that $\kappa_n := \inf_{|u|\le\ell_n} f_U(u) > 0$ with
$\kappa_n^{-1} = o(n^c)$ for some $c > 0$.
\end{assumption}

\begin{assumption}[Bandwidth conditions]\label{ass:app5}
$h_u, h_w, h_x \to 0$ with $(n h_x^p h_u^3)^{-1} \ln n \to 0$.
\end{assumption}

\begin{assumption}[LF qualification]\label{ass:app6}
There exist $M_\eta < \infty$ and $\eta > 0$ such that, for the LF spectral
filter, $\sup_{N^{-1} \le t \le \|\mathsf T'_{b_\dagger}\|^2}
t^{\beta/2 - \eta} \le M_\eta N^{-\beta/2} N^\eta$.
\end{assumption}

% ----------------------------------------------------------------------------
% PATCH-CR1: Sub-Gaussian concentration assumption.
% This is the assumption that replaces the (unsupported) appeal to ass:app4
% for sub-exponential tails in the original proof.
% ----------------------------------------------------------------------------

\begin{assumption}[Sub-Gaussian conditional response]\label{ass:subgauss}
There exists $\sigma_Y^2 < \infty$ such that, for all $(x,w)$ and all
$t \in \IR$,
\[
E\!\bigl[\exp\bigl(t\,(Y - E[Y \mid X = x, W = w])\bigr) \,\big|\, X = x, W = w\bigr]
\;\le\; \exp\!\bigl(\tfrac12\,\sigma_Y^2\, t^2\bigr).
\]
Equivalently, $U \mid X, W$ is sub-Gaussian with variance proxy
$\sigma_Y^2 / R_\sigma^2$ (this equivalence uses Assumption
\ref{ass:unifbd}\textup{(iii)}).
\end{assumption}

% ----------------------------------------------------------------------------
% PATCH-CR2: A-priori boundedness of iterates in derivative coordinates.
% Two equivalent options were on the table; we adopt the moment-bound version
% (uniformly bounded $L^4$ norm of iterates) because it is the cleanest to
% propagate in the off-event bookkeeping.
% ----------------------------------------------------------------------------

\begin{assumption}[Uniform $L^4$ bound on iterates]\label{ass:iterbd}
There exists a finite constant $C_4 \ge 1$ such that, for all $n$ large enough,
\[
\sup_{0 \le m \le M_n} E\,\bigl\|\hat b'_m\bigr\|^4 \;\le\; C_4.
\]
% PATCH-R4-B4: derivation chain remark.
Under Assumption~\ref{ass:rateconv}\textup{(i)} at $l_0 \ge 4$ and
Assumption~\ref{ass:unifbd}, the $L^4$ envelope $C_4$ is finite and
bandwidth-free; the assumption packages the resulting uniform bound
without re-deriving the chain.
\end{assumption}

% ----------------------------------------------------------------------------
% PATCH-MN9: Growth of the maximal Lepskii index.
% ----------------------------------------------------------------------------

% PATCH-R3-G2-d: tightened M_n growth from delta_n^{-2} to delta_n^{-1}.
\begin{assumption}[Maximal Lepskii index]\label{ass:mn-growth}
The Lepskii horizon satisfies $M_n \to \infty$ with $M_n = O(\delta_n^{-1})$.
In particular, $m^*(n) \le M_n$ eventually whenever the source-condition
exponent satisfies $\beta_\epsilon \ge 1$, where
$m^*(n) := \arg\min_{1 \le m \le M_n}[B(m) + \tilde V(m)]$ is the oracle
balancing index defined in Proposition \textup{\ref{prp:lepskii}}.
\end{assumption}

\begin{remark}[Scope trade-off implied by the tightened $M_n$]\label{rem:scope}
% PATCH-R3-G2-e: spell out the adaptivity scope restriction.
Round-2 of this proof attempted $M_n \asymp \delta_n^{-2}$. The rate
algebra in the propagator-residual absorption step (Step~3 of
Proposition~\ref{prp:lepskii}) does \emph{not} close at that horizon
even under the strengthened
Assumption~\ref{ass:rateconv}\textup{(iv)}: a quadratic-remainder
argument (replacing the first-order Lipschitz step \eqref{eq:Tlip-historical} by the
$C^2$-Fr\'echet hypothesis of
Assumption~\ref{ass:frediff}\textup{(ii)}) requires
$\delta_n^{2} m^{2} V(m)^{2} = o(1)$ at $m = M_n$; with
$V(M_n)^2 \asymp \delta_n^{2}M_n$, this forces $M_n = O(\delta_n^{-1})$.
The cost of this tightening is that the calibrated optimal index
$m^{*}(n) \asymp \delta_n^{-2/(\beta_\epsilon + 1)}$ of
Corollary~\ref{cor:rates} satisfies $m^{*}(n) \le M_n$ only when
$\beta_\epsilon \ge 1$. For source-condition exponents
$\beta_\epsilon < 1$ (very rough source conditions), the optimal index
exceeds $M_n$ and the Lepskii rule is \emph{not adaptive} under the
tightened Assumption~\ref{ass:mn-growth}; instead, the rule selects
$\hat m = M_n$ and the realised rate becomes $V(M_n)$, strictly slower
than the oracle. Lepskii adaptation under
Proposition~\ref{prp:lepskii} therefore covers the regime
$\beta_\epsilon \ge 1$ only --- a substantive narrowing of the result's
scope relative to the source paper's stated theorem.
\end{remark}
